\subsection{* The case of spaces of bounded continuous functions}

For every topological space X, let \(\mathcal{C}^{b}(X)\) denote the Banach space of all continuous and bounded mappings from X into K, with the norm defined by

\[
\| f \| = \sup _ {x \in X} | f (x) |\tag{6}
\]

(GT, X, § 3, No.~2). When X is compact, every continuous function on X is bounded (GT, IV, § 6, No.~1), and we write \(\mathcal{C}(X)\) for \(\mathcal{C}^b(X)\) .

In this and the following section, we shall use the following lemma, which is a particular case of Lebesgue's theorem (INT, IV, 2nd ed.~§ 4, No.~3, th. 2) on account of the interpretation of the elements of \(\mathcal{C}(\mathrm{X})'\) as measures on X.

Lemma 2. --- Let X be a compact space. If a sequence \((f_n)_{n\in \mathbf{N}}\) is bounded in \(\mathcal{C}(\mathbf{X})\) and converges simply on X to a continuous function \(f\) , then \(\mu(f) = \lim_{n\to \infty}\mu(f_n)\) for every \(\mu\) in \(\mathcal{C}(\mathbf{X})'\) .

PROPOSITION 3. --- Let X be a compact space, and let A be a bounded subset of \(\mathcal{C}(\mathrm{X})\) . For A to be relatively compact for the topology of simple convergence, it is necessary and sufficient that it is relatively compact for \(\sigma(\mathcal{C}(\mathrm{X}), \mathcal{C}(\mathrm{X})')\) .

The topology of simple convergence is Hausdorff and coarser than \(\sigma(\mathcal{C}(X), \mathcal{C}(X)')\) , hence the condition stated is sufficient (GT, I, § 9, No.~4, cor. 3).

Now suppose that A is relatively compact for the topology of simple convergence. Let \((f_{n})_{n\in\mathbf{N}}\) be a sequence of elements of A. By prop. 2 (IV, p.~33), there exists a sequence \((f_{n_{k}})\) extracted from \((f_{n})\) and converging simply to a continuous function f.~By lemma 2, the bounded sequence \((f_{n_{k}})\) tends to f for \(\sigma(\mathcal{C}(\mathrm{X}),\mathcal{C}(\mathrm{X})')\) . Then Šmulian's theorem (IV, p.~36, th. 2) shows that A is relatively compact for \(\sigma(\mathcal{C}(\mathrm{X}),\mathcal{C}(\mathrm{X})')\) .

COROLLARY. --- Let S be a topological space and A a bounded subset of \(\mathcal{C}^{b}(S)\) . The following conditions are equivalent :

\begin{enumerate}
\def\labelenumi{(\roman{enumi})}
\item
  A is relatively compact for \(\sigma(\mathcal{C}^b(\mathbf{S}),\mathcal{C}^b(\mathbf{S})')\) ;
\item
  if \((f_n)_{n\in \mathbb{N}}\) is a sequence of elements of \(\mathbf{A}\) and \((x_m)_{m\in \mathbb{N}}\) is a sequence of points of \(\mathbf{S}\) such that the iterated limits
\end{enumerate}

\[
\gamma = \lim _ {m \rightarrow \infty} \lim _ {n \rightarrow \infty} f _ {n} (x _ {m}), \quad \delta = \lim _ {n \rightarrow \infty} \lim _ {m \rightarrow \infty} f _ {n} (x _ {m})
\]

exist, then \(\gamma = \delta\) .

Let X be the Stone-Čech compactification of S (GT, IX, § 1, No.~6) and \(\alpha\) the canonical mapping from S into X. Put \(D = \alpha(S)\) . The mapping \(\phi: f \mapsto f \circ \alpha\) is an isomorphism from the normed space \(\mathcal{C}(X)\) onto the normed space \(\mathcal{C}^b(S)\) ; put \(\tilde{A} = \phi^{-1}(A)\) . Since X is compact and D is dense in X, the prop. 2 (IV, p.~33) shows that condition (ii) is equivalent to the compactness of \(\tilde{A}\) for the topology of simple convergence. The equivalence of (i) and (ii) then follows from prop. 3. *

